{\large\textbf{Earth's magnetic field measurement (10 points)}}

\textbf{Introduction}

This problem aims to measure the horizontal component of the Earth's magnetic
field. A magnet will first be characterized using a so called Gouy balance,
before being used to measure this magnetic field.

In the entire problem, uncertainties are expected to be determined only from
the fits and not from the individual experimental points.

\textbf{Equipment list}

\begin{center}\includegraphics[width=0.99\linewidth]{figures/IPhO_2025_Q4_fig1.jpg}\end{center}

\begin{center}\textbf{Fig. 1.} Photographs of all equipment.\end{center}

The list of equipment is given below and illustrated in Fig. 1. The number of
items is indicated between [] when it is greater than one. Students should ask
for help if something appears not to be working.
\begin{itemize}
\item \textbf{(a)} Magnets [3]. One magnet is attached to the force sensor (b)
and should not be removed. Another magnet is inserted into the pod (f) and
should not be removed until specified. The last one will be used in A.5. All
magnets are supposed identical.
\item \textbf{(b)} Force sensor. Connected to the Arduino (c), this sensor
measures the force along its axis, noted $m_\mathrm{f}$, in grams-force (\textquotedbl g\textquotedbl),
which is the force experienced by a 1-gram mass on the earth's surface in the
gravity field ($g_0 = 9.81\,\mathrm{m\cdot s^{-2}}$) . One of the magnets (a)
is attached to it. Each time it is switched back on, the sensor display is
reset to 0, regardless of the situation. \textit{This sensor must not be
subjected to forces in excess of 200 grams. It needs to be unpacked
carefully.}
\item \textbf{(c)} Arduino with digital display. This element is used to power
the coils (e) and to perform force and magnetic field measurements, displayed
directly in gram-force (\textquotedbl g\textquotedbl) and mT. The battery (j) powering the Arduino must
be connected to slot (i), and the battery (j) powering the coils (e) to slot
(ii) (pay attention to connection polarity). The force sensor (b) and magnetic
field sensor (d) should be connected to slots (iv) and (iii) respectively, and
the coil power cables to slots (v). A switch (vi) closes the coil supply
circuit (indicated by an LED), whose electric current can be controlled in
(vii).
\item \textbf{(d)} Magnetic field sensor with ruler. Connected to the Arduino
(c), this probe measures the field $B_z$ along the direction
$\overrightarrow{e_z}$ of the ruler, in mT.
\item \textbf{(e)} Coils in anti-Helmholtz configuration (wound in opposite
directions). These coils must be connected in series with the ammeter (g) and
to the Arduino (c) to create a magnetic field.
\item \textbf{(f)} Metallic stand on a wooden base, with suspended pod where a
magnet (a) is initially inserted, and with angle markers. The detailed assembly
of this device is explained below.
\item \textbf{(g)} Multimeter. Only used as an ammeter at the
$10\,\mathrm{A}$ range. If left inactive, the multimeter switches off, and must
be switched back on by returning it to the “OFF” position. Do not use the two
cables supplied in the multimeter case.
\item \textbf{(h)} Electric wires [3].
\item \textbf{(i)} $40\,\mathrm{cm}$ ruler.
\item \textbf{(j)} $9\,\mathrm{V}$ batteries [3]. Their capacity is of the
order of $300\,\mathrm{mA\cdot h}$.
\item \textbf{(k)} Chronometer.
\item \textbf{(l)} Adhesive paste. Can be used for the entire problem.
\end{itemize}

\begin{center}\includegraphics[width=1.0\linewidth]{figures/IPhO_2025_Q4_fig2.jpg}\end{center}

\begin{center}\textbf{Fig. 2.} Use of sensors inside the anti-Helmholtz coils.\end{center}

\textbf{Use of sensors interfaced with the Arduino (Fig. 2)}

The magnetic field sensor (d) can slide in the coils (e) as shown in (i), while
measuring the field on their axis. The $z = 0$ position for the sensor is shown
in (ii), and $z$ increases as it moves inside the coils.

The force sensor (b) is inserted into the coils as shown in (iii), before
turning the coil as in (iv) so that the transducer is vertical. \textit{To do
this, be sure to route the electrical wires through the gutters provided.}

\textbf{Installation of equipment (f) (Fig. 3), to be mounted only before
starting part B, with a $34\,\mathrm{cm}$ wire}
\begin{itemize}
\item Insert the metal post (f0a) into the wooden plate with plastic feet (f0b)
to form the stand (f0).
\item The part (f1) is located on the lower part and marks the angle of the
pod. Install the arm (f1b) on the metal post by means of a screw (f4), then fix
the part (f1a) on it with a second screw (f4).
\item The part (f2) is located on the upper part and hold the wire supporting
the pod. Install the arm (f2b) on the metal post by means of a screw (f4), then
insert the part (f2a) on it.
\item To build the pod (f3), insert the inertia bar (f3b) and a toothpick (f3c)
into the carrier part (f3a) on which a magnet (a) is already inserted. Insert
the wire supporting the pod into the part (f2a), and secure it with a screw
(f4). Turning part (f2a) changes the angle at which the wire is attached. The
toothpick allows to precisely measure the angular position of the pod.
\end{itemize}

\begin{center}\includegraphics[width=0.95\linewidth]{figures/IPhO_2025_Q4_fig3.jpg}\end{center}

\textbf{Fig. 3.} Installation of the pod on the metallic stand. Parts (f1a),
(f1b), (f2a), (f2b), and (f3a) are shown from two different angles. There are
four identical (f4) plastic screws.

\textbf{Part A. Gouy balance and magnetic moment}

\textbf{Modeling}

We assume that a magnet can be treated as a magnetic dipole of magnetic moment
$\overrightarrow{m}_\mathrm{m}$. The force experienced by such a dipole of
magnetic moment $\overrightarrow{m}_\mathrm{m} = m_\mathrm{m}\overrightarrow{e_z}$
in a magnetic field $\overrightarrow{B} = B(z)\overrightarrow{e_z}$ is
\begin{equation}
\overrightarrow{F}(z) = m_\mathrm{m}\frac{dB(z)}{dz}\overrightarrow{e_z}\,. \tag{1}
\end{equation}
When an electric current $i$ flows through the anti-Helmholtz coils, the field
$\overrightarrow{B}$ along the unit vector $\overrightarrow{e_z}$ of revolution
axis is
\begin{equation}
\overrightarrow{B}(z) = \alpha\, i(z - z_0)\overrightarrow{e_z}\,. \tag{2}
\end{equation}
This equation is only valid near the center of the device, denoted by
$z = z_0$.

\textbf{Magnetic field in the coils}

\textbf{A.1}\quad Estimate numerically the typical operating time $\tau$ of one
of the batteries used in the experiment, with an electric current of the order
of $2\,\mathrm{A}$. \hfill 0.2pt

\textit{This result must be taken into account when developing the protocols
later on, knowing that the coils are only used in part A. Note that a spare
battery is available if required.}

Insert the magnetic field sensor into the coils, as shown in Fig 2. See also
this figure for the identification of the sensor position in the coils.

\textbf{A.2}\quad At a fixed electric current $i_0 \simeq 1.0\,\mathrm{A}$,
measure and plot the magnetic field $B_z$ as a function of the position $z$ of
the sensor on the axis of the coils. Identify the largest region
$[z_\text{min}, z_\text{max}]$ where the magnetic field is experimentally
linear with respect to position. \hfill 0.8pt

\textbf{A.3}\quad By placing the sensor at two positions $(z_1, z_2)$ in this
region of linear dependency, draw a curve to verify the electric current
dependency of $\overrightarrow{B}$ given by equation (2), and determine the
value of $\alpha$, with its uncertainty. \hfill 0.9pt

\textbf{Gouy balance}

Remove the magnetic field sensor from the coils, and carefully place the force
sensor inside, as described in Fig. 2, with particular attention to the
placement of electrical wires in the gutters.

\textbf{A.4}\quad Perform experimental measurements of the gram-force
$m_\mathrm{f}$ as a function of current $i$. Draw an appropriate plot to
determine the value of the magnetic moment $m_\mathrm{m}$ of the magnet, with
its uncertainty. \hfill 0.8pt

\textbf{Alternative measurement of the magnetic moment}

In the dipolar approximation, the magnetic field of a magnet of magnetic moment
$m_\mathrm{m}$ on its revolution axis $z$ is
\begin{equation}
B_z(z) = \frac{\mu_0 m_\mathrm{m}}{2\pi(z - z_\mathrm{a})^3}, \tag{3}
\end{equation}
where $z_\mathrm{a}$ is not necessarily the geometric center of the magnet, and
where $\mu_0 = 4\pi 10^{-7}\,\mathrm{H\cdot m^{-1}}$.

\textbf{A.5}\quad Measure the magnetic field $B_z$ along the revolution axis of
the free magnet, as a function of distance $z$. Draw a curve to verify the
model given Eq. (3), showing its experimental deviations. Deduce a new value for
$m_\mathrm{m}$, with uncertainty. \hfill 1.3pt

\textbf{A.6}\quad Given the two results obtained in A.4 and A.5, propose a final
experimental value of $m_\mathrm{m}$ with its uncertainty. \hfill 0.2pt

\textbf{Part B. Determining the earth's magnetic field}

\textbf{Modeling}

We now study the oscillating motion of the magnet in a horizontal plane to
estimate the value of the horizontal component $B_\mathrm{e}$ of the Earth's
magnetic field, see Fig. 3 and the assembly instructions above Fig.3. The pod
(f3), containing the magnet, is subjected to two torques around the vertical
axis:
\begin{itemize}
\item the torque of the wire, modeled as
$\varGamma_\mathrm{f} = -\frac{C_\mathrm{f}}{L}(\theta - \theta_0)$, where
$C_\mathrm{f}$ is a constant and $L$ the total length between the two
attachments of the wire, and $\theta_0$ corresponds to the angle for which the
wire is not twisted,
\item the torque of the Earth’s magnetic fields, given by
$\varGamma_\mathrm{e} = -m_\mathrm{m}B_\mathrm{e}\sin(\theta - \theta_\mathrm{e})$,
when the angular position of the Earth's magnetic field is given by the angle
$\theta_\mathrm{e}$.
\end{itemize}
Denoting $J$ the unknown moment of inertia of the pod and magnet assembly
around the vertical axis, the angular momentum theorem gives
\begin{equation}
J\frac{\mathrm{d}^2\theta}{\mathrm{d}t^2} = \varGamma_f + \varGamma_e
= -\frac{C_\mathrm{f}}{L}(\theta - \theta_0)
- m_\mathrm{m}B_\mathrm{e}\sin(\theta - \theta_\mathrm{e}). \tag{4}
\end{equation}
When the $\sin(\theta - \theta_\mathrm{e}) \simeq \theta - \theta_\mathrm{e}$
approximation is valid, this leads to an sinusoidal oscillation at a period
$T$. For this part, adhesive past (l) is moldable into any shape or size and
attachable to other devices.

\textit{Caution:}\quad To avoid disturbance from external magnetic fields, the
magnet must be placed at least $20\,\mathrm{cm}$ away from any metal object or
magnetic source (including the other magnets).

\textbf{Experimental set-up and first measurement}

\textit{For questions B.1 to B.5, set the length of the wire to
$L = 34\,\mathrm{cm}$ and make sure that it is not twisted. In this setting, we
begin by assuming that the torque from the wire is negligible with respect to
the torque from the Earth's magnetic field, a hypothesis to which we will
return later.}

To align $\theta_0$ with $\theta_\mathrm{e}$, use piece (f2a) to adjust
$\theta_0$ so the pod (f3) does not rotate when the magnet is removed. Then
reinsert the magnet in the pod, and keep $\theta_0$ unchanged until question
B.5.

\textbf{B.1}\quad Propose an experimental protocol to determine
$B_\mathrm{e}$. Introduce the different quantities you will measure and their
units. Depict these quantities on a detailed schematic, and relate them to
those given in the instructions through an equation. For each quantity, specify
whether it is fixed (F) or varies (V) throughout the protocol. \hfill 0.3pt

\textbf{B.2}\quad Using the protocol described above, draw a graph to determine
a first value of $B_\mathrm{e}$, with its uncertainty. \hfill 1.1pt

\textbf{Evaluation of the torque from the wire}

\textbf{B.3}\quad Keeping $L = 34\,\mathrm{cm}$, study the motion of the pod
without the magnet, and determine the value of $C_\mathrm{f}$, with
experimental uncertainty: perform one period measurement for two system
configurations. Specify the equation relating $C_\mathrm{f}$ to the measured
quantities. \hfill 0.7pt

\textbf{B.4}\quad Using previous measurements, give the expression and
determine numerically the critical length $L_\mathrm{c}$ for which the
amplitude factors $C_\mathrm{f}/L$ and $m_\mathrm{m}B_\mathrm{e}$ of the
$\varGamma_\mathrm{f}$ and $\varGamma_\mathrm{e}$ torques are equal. In
question B.2, what was the ratio $(C_\mathrm{f}/L)/(m_\mathrm{m}B_\mathrm{e})$
? Choose from the intervals: [0\,\%, 1\,\%[ ; [1\,\%, 5\,\%[ ; [5\,\%, 20\,\%[ ;
[20\,\%, 50\,\%[ ; [50\,\%, $\infty$\%[. \hfill 0.3pt

\textbf{Static regime measurement}

We now propose a static measurement of the Earth's magnetic field. Reinsert the
magnet into the pod. Use piece (f2a) in Fig. 3 to adjust the angular position
$\theta_0$, causing the wire to twist.

\textbf{B.5}\quad Still at a fixed length of $L = 34\,\mathrm{cm}$, draw an
appropriate plot to study how the equilibrium position of the magnet
$\theta_\mathrm{eq}$ depends on the angle $\theta_0$, and determine a second
value of $B_\mathrm{e}$, with its uncertainty. \hfill 1.1pt

\textbf{B.6}\quad Vary the length $L$ and repeat the previous study for two
other lengths to verify the $L$ dependence of the wire torque. Using a final
graph that summarizes all the dependencies, determine a new value for
$B_\mathrm{e}$ , with its uncertainty. \hfill 2.3pt
